\documentstyle[%


righttag%


,11pt%


,amssymb%


,russian%               remove in Eng.


,izvru%                 Set izven% in Eng.


]{amsart}





%      Math definitions





\newcommand{\ctg}{\operatorname{ctg}}


\newcommand{\tts}[1]{}





%\hoffset=-15mm                  For Eng.





\setcounter{page}{10}


\god{1998}


\nomer{3~(430)} %                Remove in Eng.


%\nomer{Vol.\,42, No.\,2}        Set in Eng.


% \str{\pageref{begin}--\pageref{end}}   Set in Eng.


\udk{517.}


\author{\it И.С.\,Емельянова}


% NB    ^^ remove in Eng.


\title{Один случай решения задачи Гесса\\


в тригонометрических функциях}





% \pagestyle{myheadings}         Set in Eng.





\pagestyle{plain} %              Remove in Eng.


\begin{document}





% \thispagestyle{plain}          Set in Eng.





\maketitle


\label{begin}





В предлагаемой статье найдено точное аналитическое частное решение


уравнения


Гамильто\-на--Якоби для классической задачи о вращении тяжелого


твердого тела вокруг закрепленной точки. Предполагается, что


координаты центра тяжести тела, форма его трехосного эллипсоида


инерции в неподвижной точке, начальные условия и характеристика


силового поля удовлетворяют четырем условиям, два из которых


совпадают с условиями В.\,Гесса--Г.Г.\,Аппельрота [1]. Выполнение


третьего условия задачи В.\,Гесса--Г.Г.\,Аппельрота первоначально не


предполагается. Тем не менее показано, что найденное частное решение


удовлетворяет всем трем условиям задачи В.\,Гесса--Г.Г.\,Аппельрота.


Новое нетривиальное условие связывает значение константы


интеграла


энергии с параметрами тела и включает ускорение силы тяжести.


Примечательно, что форма нового налагаемого ограничения имеет


структуру, аналогичную основному условию В.\,Гесса. Как


известно [1],


В.\,Гессом и Г.Г.\,Аппельротом доказана принципиальная


возможность


получения однозначного аналитического решения задачи в указанном


ими случае. Однако само это решение не получено.


В [2] показано, что задача приводит к решению


обыкновенного


дифференциального уравнения типа Риккати с двоякопериодическими


коэффициентами. В свою очередь коэффициенты определяются неявно


через $P$-функции Вейертшрасса. Приведенное в статье новое решение


имеет достаточно простую структуру и содержит только


тригонометрические функции.





Функция Гамильтона $H$ тяжелого твердого тела, вращающегося


вокруг неподвижной точки, может быть представлена в следующей


форме [3], [4]:


\begin{multline}


H=\{A[P^2_\theta+(P_\psi-P_\varphi\cos\theta)^2


/\sin^2\theta]+(B-A)[P_\theta\cos\varphi\sin\theta+ \\


+(P_\psi-P_\varphi\cos\theta)\sin\varphi]^2/sin^2\theta\}/2AB+


P^2_\varphi/2C-mg[(x\sin\varphi+y\cos\varphi)


\sin\theta+z\cos\theta],


\end{multline}


где $A$, $B$, $C$ --- главные осевые моменты инерции в точке


закрепления


тела; $\varphi$, $\psi$, $\theta$ --- углы Эйлера (углы собственного


вращения, прецессии и


нутации); $P_\varphi$, $P_\psi$, $P_\theta$ --- обобщенные импульсы;


$x$,


$y$, $z$ ---координаты


центра тяжести в системе координат, начало которой находится в


неподвижной точке, а оси совпадают с главными осями инерции в этой


точке; $m$ --- масса тела, $g$ --- ускорение силы тяжести.





Функция Гамильтона (1) приобретает более компактный вид после


введения безразмерных переменных и параметров ([5], [6], с.\,99):


\begin{multline}


H=\alpha p^2_\varphi/2+\beta[p^2_\theta+(p_\psi-p_\varphi


\cos\theta)^2/\sin^2\theta]/2+(1-\beta)[p_\theta\cos\varphi\sin\theta+\\


+(p_\psi-p_\varphi\cos\theta)\sin\varphi]^2/2\sin^2\theta-


(x_0\sin\varphi+y_0\cos\varphi)\sin\theta-z_0\cos\theta


\end{multline}


($\alpha=A/C$; $\beta=A/B$ --- параметры, определяющие форму эллипсоида


инерции; $p_\varphi$, $p_\psi$, $p_\theta$ --- безразмерные обобщенные


импульсы;


$x_0$, $y_0$, $z_0$ --- безразмерные проекции радиуса-вектора центра


тяжести на выбранные координатные оси).





Уравнение Гамильтона--Якоби [6], [7], соответствующее безразмерной


функции Гамильтона (2), запишется следующим образом:


\begin{multline}


\partial w/\partial t+\alpha(\partial w/\partial\varphi)^2/2


+\beta[(\partial w/\partial\theta)^2+


(\partial w/\partial\psi-\partial w/\partial\varphi\cos\theta)^2


/\sin^2\theta]/2+(1-\beta)\{(\partial w/\partial\theta)


\cos\varphi\sin\theta+\\


%


+[\partial w/\partial\psi-(\partial w/\partial\varphi)


\cos\theta]\sin\varphi\}^2/2\sin^2\theta-


(x_0\sin\varphi+y_0\cos\varphi)\sin\theta-z_0\cos\theta=0,


\end{multline}


где $w$ --- искомая функция уравнения Гамильтона--Якоби.





Поскольку переменные $t$ и $\psi$ не содержатся явно в функции


Гамильтона $H$ (2) и, следовательно, в уравнении (3), то полное решение


уравнения (3) может быть найдено в форме, содержащей эти переменные


аддитивно:


\begin{equation}


w=-E_0t+p_0\psi+w_1(\varphi,\theta).


\end{equation}


Константы $E_0$ и $p_0$ определяют значения полной механической


энергии тела и момента количества движения в проекции на вертикаль


(система допускает интегралы энергии и площадей).





Подставляя (4) в уравнение (3), имеем


\begin{multline}


\alpha(\partial w_1/\partial\varphi)^2+\beta\{(


\partial w_1/\partial\theta)^2+[p_0-


(\partial w_1/\partial\varphi)\cos\theta]^2/\sin^2\theta\}+


(1-\beta)\{(\partial w_1/\partial\theta)\cos\varphi\sin\theta+\\


%


+[p_0-(\partial w_1/\partial\varphi)\cos\theta]\sin\varphi


\}^2/\sin^2\theta=2[E_0+(x_0\sin\varphi+y_0\cos\varphi)


\sin\theta+z_0\cos\theta].


\end{multline}


Если удается построить решение


$w_1(\varphi,\theta,E_0,p_0,a)$ уравнения (5), в


которое входит дополнительная константа интегрирования $a$, и тем


самым полное решение


$w(\varphi,\theta,\psi,t,E_0,p_0,a)$ уравнения (3)


\begin{equation}


w=-E_0t+p_0\psi+w_1(\varphi,\theta,E_0,p_0,a),


\end{equation}


то с помощью теоремы Гамильтона--Якоби [7] можно построить полное


решение


\begin{equation}


\begin{split}


&\partial w/\partial\varphi=p_\varphi;\quad


\partial w/\partial\psi=p_\psi;\quad


\partial w/\partial\theta=p_\theta;\\


&\partial w/\partial a=b_1;\quad


\partial w/\partial E_0=b_2;\quad


\partial w/\partial p_0=b_3


\end{split}


\end{equation}


системы канонических уравнений


\begin{equation}


\begin{split}


&d\varphi/dt=\partial H/\partial p_\varphi;\quad


d\psi/dt=\partial H/\partial p_\psi;\quad


d\theta/dt=\partial H/\partial p_\theta;\\


&dp_\varphi/dt=-\partial H/\partial\varphi;\quad


dp_\psi/dt=-\partial H/\partial\psi;\quad


dp_\theta/dt=-\partial H/\partial\theta,


\end{split}


\end{equation}


в которой в качестве $H$ выступает функция (2), а величины


$b_1$, $b_2$, $b_3$ в решении (7) --- константы.





Известно, что рассматриваемая задача полностью интегрируется в


случаях, обнаруженных Л.\,Эйлером ($x_0=y_0=z_0=0$), Ж.\,Лагранжем


($x_0=y_0=0$, $\beta=1$) и С.В.\,Ковалевской ($z_0=0$, $\alpha=2$,


$\beta=1$). В этих случаях


находится функция $w$ (6), содержащая константы $E_0$, $p_0$ и $a$. Во


всех трех случаях 11 параметров


\begin{equation}


\alpha,\; \beta, \; x_0, \; y_0, \; z_0, \; E_0, \;


p_0, \; a, \; b_1, \; b_2, \; b_3


\end{equation}


связаны тремя условиями, и 8 параметров могут быть выбраны


произвольно. Поскольку константы $b_2$ и $b_3$ определяют выбор начала


отсчета циклической координаты $\psi$ и времени $t$ и, следовательно,


не являются существенными, можно считать, что случай Л.\,Эйлера


допускает 6 существенных безразмерных параметров, характеризующих


форму эллипсоида инерции в неподвижной точке и начальные условия.


Что касается случаев Ж.\,Лагранжа и С.В.\,Ковалевской, то в связи с тем,


что эллипсоид инерции в неподвижной точке имеет круговое главное


сечение ($\beta=1$, т.е. $A=B$), координата $\psi$ становится циклической,


и константа $b_1$, входящая в число параметров (9), также перестает быть


существенной. Таким образом, случаи Ж.\,Лагранжа и С.В.\,Ковалевской


допускают 5 существенных безразмерных параметров. При наложении


условия кинетической симметрии ($\alpha=\beta=1$) без ограничения


общности


можно положить: $x_0=y_0=0$, и число существенных безразмерных


параметров снижается до четырех.





Продемонстрируем на наиболее известных случаях существования


полной или частной интегрируемости задачи (8), найденных к началу


нашего столетия, каким числом существенных безразмерных параметров


можно характеризовать полученные решения (таблица). Анализ


таблицы показывает, что среди случаев существования частных


интегралов наибольшим числом свободных параметров обладают случаи


В.\,Гесса--Г.Г.\,Аппельрота, О.\,Штауде и


Д.Н.\,Бобылева--В.А.\,Стеклова--Б.К.Млодзеевского.


\medskip





{\hspace{12cm} Таблица}


\smallskip





\begin{tabular}{|lccc|}


\hline


Автор случая (или его описание) & Число & Число & Дата опубл.\\


\ & огранич. & своб. парам. & \ \\


\hline


Эйлер Л. & 3 & 6 & 1758 \\


Лагранж Ж. & 3 & 5 & 1773 \\


Кинетич. симметрия & 2 & 4 & 1773 \\


Ковалевская С.В. & 3 & 5 & 1886 \\


Делоне Н.Б. & 4 & 4 & 1892 \\


Гесс В., Аппельрот Г.Г. & 3 & 5 & 1890 \\


Штауде О. & 3 & 5 & 1894 \\


Бобылев Д.Н., Стеклов В.А. & 3 & 5 & 1894 \\


Млодзеевский Б.К., Бобылев Д.Н., Стеклов В.А. & 5 & 3 & 1896 \\


Стеклов В.А. & 4 & 4 & 1899 \\


Горячев Д.Н. & 5 & 3 & 1899 \\


Горячев Д.Н., Чаплыгин С.А. & 4 & 3 & 1900 \\


Чаплыгин С.А. & 5 & 3 & 1904 \\


\hline


\end{tabular}


\medskip





Следует отметить, что число существенных параметров является


лишь одной из многочисленных характеристик случаев интегрируемости


задачи (8). Немаловажную роль в понимании найденных решений играет


простота получаемого аналитического выражения, а также


геометрической, топологической или иной трактовки результата. До


настоящего времени классическая задача (8) для функции Гамильтона (2)


оставляет много нерешенных вопросов и продолжает привлекать


постоянное внимание математиков и механиков.





Остановимся на новом полученном нами частном решении задачи


(8), представленной в форме уравнения Гамильтона--Якоби (5). Пусть


выполнены следующие четыре условия:





1) центр масс тела расположен в главной плоскости $yOz$


\begin{equation}


x_0=0;


\end{equation}





2) безразмерные координаты центра тяжести тела и параметры


трехосного эллипсоида инерции в неподвижной точке связаны условием


\begin{equation}


(\alpha-1)y^2_0=(1-\beta)z^2_0;


\end{equation}





3) константа $p_0$ в интеграле площадей полагается равной нулю


\begin{equation}


p_0=0;


\end{equation}





4) значение полной механической энергии тела связано с


положением центра тяжести и формой эллипсоида инерции условием


\begin{equation}


(\alpha-1)E^2_0=(\alpha-\beta)z^2_0.


\end{equation}





Будучи представленным в исходных размерных переменных,


условие (11) представляет одно из условий задачи


В.\,Гесса--Г.Г.\,Аппельрота


$$


B(A-C)y^2_0=C(B-A)z^2_0,


$$


а условие (13)


$$


B(A-C)E^2_0=A(B-C)(mg\,z_0)^2


$$


приобретает форму, явно демонстрирующую характер ограничения.


Последнее связывает значения главных моментов инерции тела в


неподвижной точке, положение центра масс, ускорение силы тяжести и


значение полной механической энергии тела.





Убедимся, что при выполнении условий (10)--(13) уравнение


Гамильтона--Якоби (5) допускает частное решение


\begin{equation}


w_1=2^{3/2}(E_0-y_0\cos\varphi\sin\theta-


z_0\cos\theta)^{1/2}.


\end{equation}





Подставим выражение (14) в уравнение (5) и покажем, что


имеет место тождество


\begin{multline}


\alpha y^2_0\sin^2\varphi\sin^2\theta+\beta(y^2_0\cos^2\theta


+z^2_0\sin^2\theta-y_0z_0\sin 2\theta\cos\varphi)+\\


+(1-\beta)(z_0\cos\varphi\sin\theta-y_0\cos\theta)^2


-E^2_0+(y_0\cos\varphi\sin\theta+


z_0\cos\theta)^2\equiv 0.


\end{multline}


Для доказательства справедливости тождества (15) выделим


слагаемые, содержащие в качестве множителей


$y^2_0$, $z^2_0$, $y_0z_0$ и $E^2_0$, и


преобразуем полученное выражение, умножив и поделив его на


$(\alpha-1)$ и


учитывая условия (11) и (13)


\begin{multline*}


y^2_0(\cos^2\theta+\cos^2\varphi\sin^2\theta+


\alpha\sin^2\varphi\sin^2\theta)+z^2_0


(\cos^2\theta+\cos^2\varphi\sin^2\theta+\beta\sin^2\varphi


\sin^2\theta)-E^2_0=\\


=z^2_0[(\alpha-1)(\cos^2\theta+\cos^2\varphi\sin^2\theta+\beta


\sin^2\varphi\sin^2\theta)+\\


+(1-\beta)


(\cos^2\theta+\cos^2\varphi\sin^2\theta+


\alpha\sin^2\varphi\sin^2\theta)-\alpha+\beta]/(\alpha-1)\equiv 0.


\end{multline*}





Легко убедиться, что тождество верно: достаточно выделить


множители при $\alpha$, $\beta$, $\alpha\beta$ и при свободном члене.


Все они равны нулю.


Таким образом, выражение (14) действительно является частным


решением уравнения Гамильтона--Якоби (5) при наложении ограничений


(10)-(13).





Частный интеграл (14) имеет физический смысл при положительном


значении выражения, стоящего под знаком радикала, т.е. при


соблюдении неравенства


$$


E_0>y_0\cos\varphi\sin\theta+z_0\cos\theta.


$$





Кроме того, следует потребовать, чтобы эллипсоид инерции в


неподвижной точке имел несовпадающие главные моменты инерции. Без


ограничения общности достаточно выполнить условия


\begin{equation}


\alpha>1>\beta.


\end{equation}


Для этого достаточно направить ось $Oz$ по линии наименьшего


момента инерции, и угол нутации $\theta$ будет характеризовать


отклонение от


вертикали наибольшей оси эллипсоида инерции. Отметим, что центр


масс тела с учетом неравенств (16) располагается в плоскости,


перпендикулярной средней оси эллипсоида инерции.





Обобщенные импульсы $p_\varphi$ и $p_\theta$, отвечающие решению (14),


вычисляются как частные производные от $w_1$ по соответствующим


обобщенным координатами


\begin{equation}


\begin{split}


p_\varphi&=2^{1/2}(y_0\sin\varphi\sin\theta)/


(E_0-y_0\cos\varphi\sin\theta-z_0\cos\theta)^{1/2};\\


p_\theta&=2^{1/2}(z_0\sin\theta-y_0\cos\varphi\cos\theta)


/(E_0-y_0\cos\varphi\sin\theta-z_0\cos\theta)^{1/2}.


\end{split}


\end{equation}


Выражения (17) представляют собой линейные инвариантные


соотношения (частные интегралы).





Решение (14) задачи (5) и следующие из него инвариантные


соотношения (17) получены автором с привлечением техники группового


анализа уравнения Гамильтона--Якоби (5). Построены два локальных


преобразования симметрии. Каждое из них означает, что одно из условий


(17) является инвариантом группы Ли на уровне выполнения другого


условия. В связи с громоздкостью проведенного анализа ограничимся


ссылками на источник идеи получения решения


([8]; [6], cc.\,65, 96, 104).





Нередко при описании частных случаев интегрируемости задачи о


вращении тяжелого твердого тела авторы завершают исследование


указанием соответствующих инвариантных соотношений, поскольку это


доказывает принципиальную возможность нахождения закона движения


тела [11]. Мы проанализируем найденное решение более подробно.





Покажем, что решение (14) задачи (5) соответствует частному


случаю задачи В.\,Гесса и Г.Г.\,Аппельрота


([1]; [2], c.\,161--274; [9], c.\,182--184; [10], c.\,264--272; [11],


c.\,119--124).


Найденный В.\,Гессом и Г.Г.\,Аппельротом


случай существования частного решения задачи о вращении тяжелого


твердого тела обычно представляют в переменных Эйлера


$\varphi$, $\psi$, $\theta$, $p$, $q$, $r$


(углы Эйлера и проекции угловой скорости тела на главные оси инерции


в неподвижной точке). Требуется выполнение следующих трех условий.





а) Центр масс тела расположен в главной плоскости эллипсоида


инерции в неподвижной точке; без ограничения общности можно считать


выполненным условие (10)


$$


x_0=0.


$$





б) Справедливо условие В.\,Гесса (с точностью до принятого выбора


осей)


\begin{equation}


B(A-C)y^2_0=C(B-A)z^2_0.


\end{equation}


После перехода к безразмерным величинам условие (18) совпадает с (11).





в) Наложено условие [10]


$$


Bqy_0+Crz_0=0,


$$


означающее, что проекция угловой скорости на направление


радиуса-вектора центра масс тела равна нулю. В безразмерных переменных


последнее условие приобретает вид


\begin{equation}


\alpha qy_0+\beta rz_0=0.


\end{equation}





Как известно, компоненты угловой скорости связаны с


каноническими переменными следующим образом:


\begin{equation}


\begin{split}


p&=p_0\cos\varphi-p_\varphi\cos\theta\sin\varphi;\\


q&=-\beta(p_\theta\sin\varphi+p_\varphi\ctg\theta\cos\varphi);


\quad r=\alpha p_\varphi.


\end{split}


\end{equation}


Подставляя в (20) значения $p_\theta$ и $p_\varphi$ (17), имеем


\begin{equation}


\begin{split}


p&=2^{1/2}(z_0\cos\varphi\sin\theta-y_0\cos\theta)/


(E_0-y_0\cos\varphi\sin\theta-z_0\cos\theta)^{1/2};\\


q&=-2^{1/2}\beta z_0\sin\varphi\sin\theta/(E_0-y_0


\cos\varphi\sin\theta-z_0\cos\theta)^{1/2};\\


r&=-2^{1/2}\alpha y_0\sin\varphi\sin\theta/


(E_0-y_0\cos\varphi\sin\theta-z_0\cos\theta)^{1/2}.


\end{split}


\end{equation}


Последние два соотношения из (21) свидетельствуют о том, что проекции


$q$ и $r$ угловой скорости пропорциональны, и выполняется условие (19),


что и требовалось показать.





Таким образом, исследуемая задача представляет собой частный случай


задачи В.\,Гесса-Г.Г.\,Аппельрота. Проведенный в [2] анализ задачи


В.\,Гесса-Г.Г.\,Аппельрота позволил получить качественную картину


движения тела. Нетрудно убедиться, что момент инерции тела относительно


оси, проходящей через неподвижную точку и центр тяжести, определяется


выражением $J=(\alpha\beta)^{-1}$ (в обозначениях В.\,Гесса $J=AC/B$).


Центр масс тела движется


как плоский математический маятник в однородном силовом поле,


увеличенном в $y^2_0+z^2_0$ раз по сравнению с полем силы тяжести


[10]. В


рассматриваемой задаче проекции момента количества движения тела на


прямую, соединяющую неподвижную точку с центром тяжести, и на вертикаль


равны нулю.





Как известно, П.А.\,Некрасовым общая задача


В.\,Гесса-Г.Г.\,Аппельрота сведена к интегрированию обыкновенного


дифференциального уравнения первого порядка с двоякопериодическими


коэффициентами. Последние, в свою очередь, определяются неявно через


$P$-функции Вейерштрасса ([10]; [11], c.\,122).


С помощью этого уравнения


определяются величины $p$, $q$, $r$. Задача не сведена к квадратурам,


хотя


принципиальная возможность получения однозначного решения доказана.


Введенное дополнительное условие (13) позволяет получить достаточно


простое аналитическое решение задачи.





Действительно, используя одно из уравнений В.\,Гесса


([10], c.\,268):


$$


\nu d\mu/dt=(\mu+E_0)d\nu/dt,


$$


($\nu=p^2+q^2/\beta+r^2/\alpha$;


$\mu=y_0\cos\varphi\sin\theta+z_0\cos\theta$),


строится явная зависимость


между углами $\varphi$ и $\theta$


\begin{gather*}


\varphi=\arccos f(\theta);\\


f(\theta)=-[(\alpha-\beta)^{1/2}+(\alpha-1)^{1/2}


\cos\theta]/(1-\beta)^{1/2}\sin\theta.


\end{gather*}


Решение существует при $|f(\theta)|\le 1$. Использование зависимости


$\varphi$ ($\theta$) позволяет завершить нахождение закона движения


тела.





Найденное частное решение задачи (3) оставляет свободными четыре


параметра задачи. Равное с описываемым случаем число ограничений и


свободных параметров (4 и 4) имеют случаи Н.Б.\,Делоне и В.А.\,Стеклова


(см. табл.). Для сравнения укажем, что в известном случае


Д.Н.\,Горячева--С.А.\,Чаплыгина при четырех ограничениях остаются всего


три свободных параметра, поскольку эллипсоид инерции в неподвижной


точке является телом вращения.





\begin{thebibliography}{99}





\bibitem{He}


Hess W. {\em \"Uber die Eulerschen Bewegungsgliechungen und \"uber eine


neue partikul\"are L\"ozung des Problems der Bewegung einesstarren


K\"orpers um einen festen Punkt} // Math. Ann. -- 1890. -- Bd.\,37. --


\No\,2. -- S.\,153--181.


\bibitem{Ne}


Некрасов П.А. {\em Аналитическое исследование одного случая движения


тяжелого твердого тела около неподвижной точки П.А.\,Некрасова}


// Тр. Моск. матем.


о-ва. -- 1896. -- Т.\,18. -- 114 с.


\bibitem{Ar}


Архангельский Ю.А. {\em Аналитическая динамика твердого тела}. -- М.:


Наука, 1977. -- 328 с.


\bibitem{Ch}


Чертков Р.И. {\em Метод Якоби в динамике твердого тела}. -- Л.:


Судпромгиз, 1960. -- 324 с.


\bibitem{E1}


Емельянова И.С. {\em Метод Якоби для вращающегося тяжелого тела}. // Методы


прикладного функционального анализа / под ред.


С.Н.\,Слугина. -- Н.\,Новгород: Изд-во Нижегород. ун-та, 1991. --


С.\,20--24.


\bibitem{E2}


Емельянова И.С. {\em Проблема ``симметрия--интегралы движения'' в


аналитической динамике}. -- Н.\,Новгород: Изд-во Нижегород. ун-та,


1992. -- 171 с.


\bibitem{Pa}


Парс Л. {\em Аналитическая динамика}. -- М.: Наука, 1971. -- 636 с.


\bibitem{Ov}


Овсянников Л.В. {\em Групповой анализ дифференциальных уравнений}. --


М.: Наука, 1978. -- 400~с.


\bibitem{Ra}


Раус Э.Дж.{\em Динамика системы твердых тел}. Т.\,2 / Под ред.


Ю.А.\,Архангельского и В.Г.\,Демина. -- М.: Наука, 1983. -- 544 с.


\bibitem{Go}


Голубев В.В. {\em Лекции по интегрированию уравнений движения тяжелого


твердого тела около неподвижной точки}. -- М.: Гостехиздат, 1953. --


287 с.


\bibitem{GK}


Горр Г.В., Кудряшова Л.В., Степанова Л.А. {\em Классические задачи


динамики твердого тела. Развитие и современное состояние}. -- Киев:


Наукова думка, 1978. -- 96 с.





\end{thebibliography}





\vskip 2cm





\post{\it Нижегородский государственный}{\it Поступила}


\post{\it университет}{23.05.1995}


\label{end}


\end{document}


